% !PRINTED 480
\weberpage{480}
Diesen Satz können wir nun auch in folgender Weise verallgemeinern:

3. \gesperrt{Wenn durch die Substitution $y = A(x)$ irgend eine Form $\varphi(x)$ in $\Phi(y)$ übergeht, wenn irgend eine zweite Form $\psi(\xi)$ durch die transponirte Substitution $\xi = A_1(\eta)$ in $\Psi(\eta)$ übergeht, so erhält man eine neue Transformation durch $A$, wenn man in $\psi(\xi)$ und $\Psi(\eta)$ die Vertauschungen macht}

\[
\begin{aligned}
&\xi_1^\alpha \xi_2^\beta \xi_3^\gamma \quad \text{mit} \quad \frac{\partial^{\alpha + \beta + \gamma} \varphi}{\partial x_1^\alpha \partial x_2^\beta \partial x_3^\gamma} \\
&\eta_1^\alpha \eta_2^\beta \eta_3^\gamma \quad \text{mit} \quad \frac{\partial^{\alpha + \beta + \gamma} \Phi}{\partial y_1^\alpha \partial y_2^\beta \partial y_3^\gamma},
\end{aligned}
\]

\gesperrt{wenn man also, wie man sich auch ausdrücken kann, in $\psi$ und $\Psi$ die Potenzen und Producte der Variablen $\xi, \eta$ durch die entsprechenden Ableitungen von $\varphi, \Phi$ nach den Variablen $x, y$ setzt.}

Aus der Compositionsregel der linearen Substitutionen ergiebt sich nun sofort der folgende Satz:

4. \gesperrt{Durchläuft $A$ eine Gruppe $G$, so durchläuft die transponirte Substitution $A_1$ eine Gruppe $G_1$. Sind $A, B$ zwei Elemente aus $G$ und $A_1, B_1$ die entsprechenden Elemente aus $G_1$, so sind $AB$ und $B_1 A_1$ entsprechende Elemente. Die Gruppen $G$ und $G_1$ werden zu einander contragredient genannt.}

Die beiden Gruppen $G$ und $G_1$ sind aber nur dann isomorph auf einander bezogen, wenn man dem $A_1$ nicht das Element $A$, sondern das Element $A^{-1}$ entsprechen lässt; denn dann entspricht $A_1 B_1$ dem Elemente $A^{-1} B^{-1} = (BA)^{-1}$.

Die Invarianten der Gruppe $G_1$ heissen Contravarianten der Gruppe $G$. Demnach sind auch die Invarianten von $G$ die Contravarianten von $G_1$.

Aus (3) ergiebt sich dann der folgende Satz:

5. \gesperrt{Wenn man in einer Contravariante von $G$ die Potenzen und Producte der Variablen durch die entsprechenden Ableitungen einer Invariante ersetzt, so erhält man wieder eine Invariante von $G$}
